%%%%%%%%% JPS abstract %%%%%%%%%%%%%%%%%%%%%%%%%%%%%%%%%%%%%%%%%%
\documentclass[12pt,a4paper,upLaTeX]{jsarticle}

%%%%%%%%% packages %%%%%%%%%%%%%%%%%%%%%%%%%%%%%%%%%%%%%%%%%%%%%%
\usepackage{graphicx} % Include figure files
\usepackage[%                           % 余白の設定
%mag=1400,%                              jarticle の場合（14ptに）
dvipdfm,truedimen,%
top=30truemm,bottom=20truemm,%
left=20truemm,right=20truemm]{geometry}
\pagestyle{empty}

%%%%%%%%% header %%%%%%%%%%%%%%%%%%%%%%%%%%%%%%%%%%%%%%%%%%%%%%%%
\begin{document}
\vspace{-5pt}
\begin{center}
{\gt \Large 講演概要集原稿の書き方 }\\[14pt]

{\gt \large 日本科技大$^A$ 帝都大理工$^B$ \\ 止田次郎$^A$, 潟川　学$^B$}\\[5pt]

{\large \bf How to Write an Abstract for the JPS Meeting}\\[5pt]

{\large \it $^A$Dept. of Phys., Teito Univ. $^B$Dept. of Phys, Nihon Univ. of Tech.}\\

{\large \bf J. Tomeda$^A$ and M. Gatagawa$^B$}
\end{center}

\vspace{10pt}
%%%%%%%%% main %%%%%%%%%%%%%%%%%%%%%%%%%%%%%%%%%%%%%%%%%%%%%%%%%%

この講演概要集原稿はtexのtemplateを用いて書かれています。\\
以下に挙げる決まりを守って頂ければ、最終判断は著者にあるので、多少の違いはあっても構いません。

\begin{itemize}
\item 大会名・講演番号・ページ数を入れるため余白の設定を変更しないでください。
\item タイトル、所属、氏名、英語表記は上記のような配置で記載してください。
\item 講演番号は記載しないでください。
\item 枠はつけてはいけません。
\item 原稿はフォントを埋め込んだpdfファイルを提出してください。
\item 1ファイルは2MB以下になるようにしてください。
\end{itemize}

以上のことを守って、講演概要の提出をお願いします。


%%%%%%%%%%%%%%%%%%%%%%%%%%%% Figure 1 %%%%%%%%%%%%%%%%%%%%%%%%%%%%%%%%%%%%%
\begin{figure}[h]
\begin{center}
\includegraphics[width=5cm]{JPS.eps}
\end{center}
\caption{日本物理学会のマーク。カラー図面が掲載できるようになった。 }
\end{figure}
%%%%%%%%%%%%%%%%%%%%%%%%%%%%%%%%%%%%%%%%%%%%%%%%%%%%%%%%%%%%%%%%%%%%%%%%%%%


\begin{description}
\item[注1：] 動画、URLへのハイパーリンクは不可とします。
\item[注2：] 概要集に掲載された原稿の著作権は日本物理学会に帰属します。
\item[注3：] 概要集アクセス権を購入者は物理学会のホームページで大会終了後1年間閲覧可能。その後は、アクセス権を持つ、持たないに関わらずJ-Stageで閲覧可能。また記録用DVDの購入も可能。
\item[注4：] 概要集原稿の提出には登録番号とパスワードを忘れずに。〆切に余裕をもってpdf原稿をwebから送信してください。
\end{description}

\end{document}